% Alpha-OSK: A Predictive On-Screen Keyboard for Motor-Impaired Users
%
% Build:  python docs/paper/build.py
% Source: docs/WHITEPAPER.md is authoritative. body.tex and abstract.tex are
%         GENERATED from it by build.py and must not be hand-edited: edit the
%         markdown and rebuild, or the next build silently discards your work.

\documentclass[11pt,a4paper]{article}

\usepackage[T1]{fontenc}
\usepackage[utf8]{inputenc}
\usepackage{lmodern}
\usepackage[margin=1in]{geometry}
\usepackage{microtype}

\usepackage{booktabs}
\usepackage{longtable}
\usepackage{array}
\usepackage{multirow}
\usepackage{graphicx}
\usepackage{float}
\usepackage{caption}
\usepackage{amsmath}
\usepackage{amssymb}
\usepackage{textcomp}
\usepackage{upquote}
\usepackage{fancyvrb}
\usepackage{calc}

\usepackage{xcolor}
\definecolor{linkblue}{HTML}{1A4F8A}
\usepackage[colorlinks=true,linkcolor=linkblue,citecolor=linkblue,urlcolor=linkblue]{hyperref}
\usepackage{url}

% Pandoc emits these for its list, code and image output. They are defined
% only inside pandoc's own template, and we supply our own preamble, so
% without them the build dies on the first figure with "undefined control
% sequence" rather than on anything to do with the content.
\providecommand{\tightlist}{\setlength{\itemsep}{0pt}\setlength{\parskip}{0pt}}
\providecommand{\passthrough}[1]{#1}
\newcommand{\euro}{\texteuro}

% \pandocbounded shrinks an oversized image to the text block instead of
% letting it run into the margin. This is pandoc's own definition.
\makeatletter
\newsavebox\pandoc@box
\providecommand*\pandocbounded[1]{%
  \sbox\pandoc@box{#1}%
  \Gscale@div\@tempa{\textheight}{\dimexpr\ht\pandoc@box+\dp\pandoc@box\relax}%
  \Gscale@div\@tempb{\linewidth}{\wd\pandoc@box}%
  \ifdim\@tempb\p@<\@tempa\p@\let\@tempa\@tempb\fi%
  \ifdim\@tempa\p@<\p@\scalebox{\@tempa}{\usebox\pandoc@box}%
  \else\usebox{\pandoc@box}%
  \fi%
}
\makeatother

% Inline code should not overflow the measure on a long identifier.
\usepackage{listings}
\lstset{
  basicstyle=\small\ttfamily,
  breaklines=true,
  breakatwhitespace=false,
  columns=fullflexible,
  frame=none,
  xleftmargin=1em,
}

% Pandoc marks an uncaptioned longtable with \LTcaptype{none} so the table
% counter does not advance. With the caption package loaded that resolves to
% a counter named "none", which does not exist, and the build dies on the
% first plain table. Defining it, printing as nothing, is the whole fix.
\newcounter{none}
\renewcommand{\thenone}{}

% This paper is full of long file paths and identifiers set in \texttt, and
% typewriter text is unhyphenated by default, so a single
% "google/osv-scanner-action@9a49..." runs 340pt into the margin. Allowing
% hyphenation inside \texttt plus a generous emergency stretch keeps every
% line inside the measure; the alternative is manual \allowbreak in prose,
% which the next edit to the markdown would undo.
\usepackage[htt]{hyphenat}
\sloppy
\emergencystretch=3em
\hbadness=10000

\captionsetup{font=small,labelfont=bf,width=.9\textwidth}
\setlength{\parskip}{0.55em}
\setlength{\parindent}{0pt}
\setcounter{tocdepth}{2}
\setcounter{secnumdepth}{3}

\title{\bfseries Alpha-OSK: A Predictive On-Screen Keyboard\\for Motor-Impaired Users}

\author{Owen Kent\\[.3em]
  \normalsize\texttt{owenpkent@gmail.com}\\
  \normalsize\url{https://github.com/owenpkent/alpha-osk}}

\date{September 2026\\[.3em]
  \normalsize Describing Alpha-OSK 1.3.0}

\begin{document}

\maketitle

\begin{abstract}
\input{abstract.tex}
\end{abstract}

\tableofcontents
\bigskip

\input{body.tex}

% ---------------------------------------------------------------------------
% References. Author-year, cited in prose in the body, so the list is
% alphabetical by first author.
% ---------------------------------------------------------------------------
\bibliographystyle{plainnat}
\begin{thebibliography}{99}

\bibitem{cleary1984}
Cleary, J.~G. and Witten, I.~H. (1984).
Data Compression Using Adaptive Coding and Partial String Matching.
\emph{IEEE Transactions on Communications}, 32(4), 396--402.

\bibitem{garbe2012}
Garbe, W. (2012).
\emph{SymSpell: 1000x faster spelling correction}.
\url{https://github.com/wolfgarbe/SymSpell}.

\bibitem{higginbotham2007}
Higginbotham, D.~J., Shane, H., Russell, S. and Caves, K. (2007).
Access to AAC: Present, past, and future.
\emph{Augmentative and Alternative Communication}, 23(3), 243--257.

\bibitem{klakow1998}
Klakow, D. (1998).
Log-linear interpolation of language models.
\emph{Proceedings of ICSLP}.

\bibitem{kristensson2004}
Kristensson, P.~O. and Zhai, S. (2004).
SHARK\textsuperscript{2}: A Large Vocabulary Shorthand Writing System for Pen-based Computers.
\emph{Proceedings of UIST}, 43--52.

\bibitem{mackenzie2002}
MacKenzie, I.~S. and Soukoreff, R.~W. (2002).
Text Entry for Mobile Computing: Models and Methods, Theory and Practice.
\emph{Human-Computer Interaction}, 17(2--3), 147--198.

\bibitem{trnka2008}
Trnka, K. and McCoy, K.~F. (2008).
Evaluating Word Prediction: Framing Keystroke Savings.
\emph{Proceedings of ACL-08: HLT, Short Papers}, 261--264.

\bibitem{vertanen2011}
Vertanen, K. and Kristensson, P.~O. (2011).
The Imagination of Crowds: Conversational AAC Language Modeling using Crowdsourcing and Large Data Sources.
\emph{Proceedings of EMNLP}, 700--711.

\bibitem{vertanen2015}
Vertanen, K., Memmi, H., Emge, J., Reyal, S. and Kristensson, P.~O. (2015).
VelociTap: Investigating Fast Mobile Text Entry using Sentence-Based Decoding of Touchscreen Keyboard Input.
\emph{Proceedings of CHI}, 659--668.

\bibitem{vescovi}
Vescovi, M.
\emph{Presage: An intelligent predictive text entry platform}.
\url{https://presage.sourceforge.io}.

\bibitem{ward2000}
Ward, D.~J., Blackwell, A.~F. and MacKay, D.~J.~C. (2000).
Dasher: A Data Entry Interface Using Continuous Gestures and Language Models.
\emph{Proceedings of UIST}, 129--137.

\bibitem{latinime}
\emph{AOSP LatinIME source}.
Android Open Source Project.
Reference implementation for trie-based dictionaries, weighted edit distance,
and letting the literal typed word compete against corrections.

\bibitem{uia}
Microsoft.
\emph{UI Automation overview}.
\url{https://learn.microsoft.com/en-us/windows/win32/winauto/}.

\bibitem{atspi}
\emph{AT-SPI 2: Accessibility Toolkit Service Provider Interface}.
GNOME Project.

\end{thebibliography}

\appendix

\section{Reproducing the evaluation}

Every figure in Section~8 comes from \texttt{scripts/bench/ksr.py} in the
source repository, which builds a cold-start engine in a fresh temporary
directory and never reads the developer's own model. The corpus and its
provenance are documented in \texttt{scripts/bench/data/README.md}.

\begin{lstlisting}
# Table 2, main results
python scripts/bench/ksr.py --corpus aac-dev
python scripts/bench/ksr.py --corpus aac-test

# Table 3, source ablation
python scripts/bench/ksr.py --corpus aac-test \
    --conditions full,no-fuzzy,legacy-fuzzy,ppm-merge

# Merge strategies
python scripts/bench/ksr.py --corpus aac-test \
    --conditions rank,rrf,linear,loglinear

# Table 5, robustness to an uncorrected pointer slip
python scripts/bench/ksr.py --corpus aac-test --mis-click \
    --conditions full,legacy-fuzzy

# Table 6, a systematically biased pointer simulated end to end
python scripts/bench/ksr.py --corpus aac-test --pointer 0.35,0.25,0.15

# Table 7, personalisation
python scripts/bench/ksr.py --corpus aac-dev --learn-half
\end{lstlisting}

The per-component design documents referenced throughout are in
\texttt{docs/architecture/}: \texttt{PPM.md}, \texttt{FUZZY\_RECOGNITION.md},
\texttt{HYBRID\_MERGING.md}, \texttt{NGRAM\_SEEDS.md}, \texttt{DICTATION.md},
\texttt{TELEMETRY.md}, \texttt{PLATFORM\_ARCHITECTURE.md} and
\texttt{COMPACT\_VIEW.md}; build and distribution detail is in
\texttt{docs/build/}.

\end{document}
